\section{Introduction}\label{sec:intro}

Biomedical vision-language models such as BiomedCLIP \citep{zhang2023biomedclip} classify images zero-shot from text prompts,
and a few labelled examples per class can adapt them further \citep{koleilat2025biomedcoop}.
With so few examples, however, the adapted model depends on which examples were drawn.
A support set that happens to be unrepresentative can produce a classifier much worse than the zero-shot one,
and in medical use the worst plausible outcome matters at least as much as the average.

Prompt-Anchored Residual Adaptation (PARA) \citep{kang2026para} addresses this directly.
It keeps the frozen prompt prediction as an anchor that no support set can change,
learns a visual prediction from the support set, and moves from the anchor towards it by a residual step $\eta(S)$.
The step is computed in closed form from the frozen support embeddings, from how far the class images lie from their text and how consistently they agree,
so that a support set which disagrees with the prompts earns a larger step.
PARA also introduces a repeated-support protocol that reports the worst $20\%$ of support draws (CVaR20) alongside the mean,
and reports state-of-the-art results on both.

This paper asks whether the support set actually steers PARA's step, and what happens when it does not.
We reimplement PARA on BiomedCLIP and study it on six MedMNIST datasets with its own repeated-support protocol, balanced accuracy and CVaR20.
Three datasets were used to develop every analysis; three others were chosen afterwards and serve as held-out tests,
including a rule whose constant was fixed before any held-out model was trained.
Our findings are:
\begin{enumerate}
  \item \textbf{The step is a fixed curve in the number of shots} (Section~\ref{sec:finding1}).
  In BiomedCLIP's frozen space, the discrepancy and agreement signals barely depend on the class:
  images sit about equally far from every class's text, and any set of images agrees almost as much as a real class.
  The step is therefore the same function of $K$ on every dataset ($0.33$, $0.60$ and $0.80$ at $K = 1$, $4$ and $16$),
  predictable from two constants of the pretrained model.
  A schedule that uses $K$ alone does better than PARA's step (significantly in 13 of 18 design settings, never worse).
  \item \textbf{On binary tasks the blend is a switch} (Section~\ref{sec:finding2}).
  Standardizing each branch's scores across two classes keeps only the sign of the prediction,
  so PARA returns the zero-shot prediction whenever $\eta < 0.5$ and the trained branch's prediction whenever $\eta > 0.5$, on every image.
  We prove this and confirm it on all 1,046 relevant evaluations.
  \item \textbf{The right step depends on the anchor's reliability relative to the trained branch} (Section~\ref{sec:finding3}).
  The best step in hindsight ranges from $0.17$ to $0.95$ across datasets and shot counts, and follows how good the anchor is compared with the trained branch.
  A rule based on the anchor's reliability alone, fixed in advance, improved on the design datasets and failed on the held-out ones.
  \item \textbf{Blending log-probabilities helps on multi-class tasks} (Section~\ref{sec:finding4}).
  With the same step and the same trained models, blending each branch's log-probabilities instead of its standardized scores
  raises mean balanced accuracy by $4.86$ points on multi-class datasets, but lowers the worst case where the anchor is good.
\end{enumerate}
From these we draw practical recommendations for evaluating adaptive blending (Section~\ref{sec:discussion}).
We do not claim that PARA is worse than other few-shot methods, and we do not reproduce its benchmark numbers;
our subject is how its step behaves and what a step of this kind would need to do.
% TODO: code and prompt-mapping release statement, once decided.
